\documentclass[a4paper,12pt]{article}
\usepackage[T2A]{fontenc}
\usepackage[utf8]{inputenc}
\usepackage[russian]{babel}
\usepackage{amsmath,amssymb,mathtools}
\renewcommand{\familydefault}{\sfdefault}
\usepackage{icomma}
\usepackage[left=24mm,right=22mm,top=23mm,bottom=24mm]{geometry}
\usepackage{microtype}
\usepackage{xcolor}
\usepackage{tikz}
\usetikzlibrary{calc,arrows.meta,positioning,shapes.geometric}
\usepackage{pgfplots}
\pgfplotsset{compat=1.18}
\usepackage{tabularx,booktabs}
\usepackage{enumitem}
\usepackage{fancyhdr}
\usepackage{parskip}
\usepackage{needspace}
\usepackage{setspace}
\usepackage[hidelinks]{hyperref}
\definecolor{accent}{HTML}{087B80}
\definecolor{darkaccent}{HTML}{15575A}
\definecolor{textgray}{HTML}{38444A}
\definecolor{muted}{HTML}{74868B}
\color{textgray}
\setstretch{1.16}
\setlist[enumerate]{leftmargin=*,itemsep=0.75em,topsep=0.5em}
\setlist[itemize]{leftmargin=*,itemsep=0.30em,topsep=0.3em}
\pagestyle{fancy}
\fancyhf{}
\fancyhead[L]{\footnotesize\color{muted}Чек-лист · Григорий Архипов}
\fancyhead[R]{\footnotesize\color{muted}14.10.26}
\fancyfoot[C]{\footnotesize\color{muted}\thepage}
\renewcommand{\headrulewidth}{0pt}
\renewcommand{\footrulewidth}{0pt}
\setlength{\headheight}{14pt}
\newcommand{\key}[1]{\textcolor{darkaccent}{\bfseries #1}}
\newcommand{\ruleline}{\par\vspace{0.25em}{\color{accent!70}\hrule height 0.6pt}\vspace{0.7em}}
\newcommand{\checkitem}[1]{\item[$\square$] #1}
\newlist{checklist}{itemize}{1}
\setlist[checklist]{label=$\square$,leftmargin=*,itemsep=0.34em}
\begin{document}
\begin{titlepage}
\thispagestyle{empty}
\begin{center}
\vspace*{36mm}
{\color{accent}\small МАТЕМАТИКА · ЕГЭ}\par
\vspace{18mm}
{\fontsize{33}{40}\selectfont\bfseries Чек-лист}\par
\vspace{11mm}
{\fontsize{20}{27}\selectfont\bfseries Текстовые задачи\\на движение}\par
\vspace{5mm}
{\large Повторение степеней и графиков}\par
\vspace{23mm}
{\color{muted}\large 14 октября 2026 года}\par
\vspace{7mm}
{\color{darkaccent}\large Григорий Архипов}\par
\vfill
{\small Лёвин Артём Александрович}\par
{\small эксклюзивно для Григория Архипова}\par
\vspace{27mm}
\end{center}
\begin{tikzpicture}[remember picture,overlay]
\draw[accent!50,line width=1.5pt] ([xshift=22mm,yshift=20mm]current page.south west)
.. controls ([xshift=59mm,yshift=32mm]current page.south west) and ([xshift=-60mm,yshift=2mm]current page.south east)
.. ([xshift=-22mm,yshift=20mm]current page.south east);
\end{tikzpicture}
\end{titlepage}

\section*{01. Степени: правила и ограничения}
\ruleline
\begin{checklist}
\checkitem{Я использую формулу \((a^m)^n=a^{mn}\) \textbf{только при допустимых значениях} основания и показателей.}
\checkitem{Для положительного \(a\) и \(p\ne0\) понимаю приём обратного показателя: \(a^p=b\Rightarrow a=b^{1/p}\).}
\checkitem{Помню, что чётная степень скрывает знак: \(t^4=16\Rightarrow t=\pm2\).}
\checkitem{При отрицательной степени учитываю \(t\ne0\), а дробные показатели проверяю на область определения.}
\checkitem{Умею разложить число на множители и представить его как степень.}
\end{checklist}

\key{Короткая памятка.} Если \(a>0\), то
\[
 a^{m/n}=\sqrt[n]{a^m},\qquad a^{-p}=\frac1{a^p},\qquad
 (a^m)^n=a^{mn}.
\]
Для нечётного \(n\) выражение \(t^{m/n}\) можно определять также при отрицательном \(t\); для чётного \(n\) в используемой школьной записи требуется \(t\ge0\). Отдельно проверяйте допустимость основания при отрицательном показателе.

\key{Пример 1.} Решим \(t^{5/3}=32\). Обратный показатель равен \(3/5\):
\[
 t=32^{3/5}=(2^5)^{3/5}=2^3=\boxed{8}.
\]
\key{Пример 2.} Решим \(t^{-3/5}=8\) при \(t\ne0\):
\[
 t=8^{-5/3}=(2^3)^{-5/3}=2^{-5}=\boxed{\tfrac1{32}}.
\]
\key{Пример 3.} Вычислим \(\sqrt[3]{59{,}319}\). Из разложения \(59319=3^3\cdot13^3\) получаем
\[
 59{,}319=\frac{3^3\cdot13^3}{10^3}
 =\left(\frac{39}{10}\right)^3,
 \qquad \sqrt[3]{59{,}319}=\boxed{3{,}9}.
\]
\key{Типичная ошибка.} Из \(t^4=16\) записать только \(t=2\). Правильно: \(t^4=16\iff |t|=2\iff t=\pm2\).

\clearpage
\section*{02. Графики: точка пересечения}
\ruleline
\begin{checklist}
\checkitem{Различаю абсциссу \(x\) и ординату \(y\).}
\checkitem{По точке \((x_0,y_0)\) гиперболы \(y=k/x\) нахожу \(k=x_0y_0\), где \(x_0\ne0\).}
\checkitem{Для прямой \(y=ax+b\) определяю \(a=\frac{y_2-y_1}{x_2-x_1}\), если \(x_2\ne x_1\), а затем \(b=y_0-ax_0\).}
\checkitem{В точке пересечения приравниваю формулы функций и исключаю \(x=0\) для гиперболы.}
\end{checklist}

\key{Пример.} Гипербола \(y=k/x\) и прямая с угловым коэффициентом \(a=5\) пересекаются в точке \(A(2;1)\). Найдём абсциссу второй точки пересечения.
\[
 k=2\cdot1=2,\quad b=1-5\cdot2=-9,
 \quad \frac2x=5x-9,\; x\ne0.
\]
\[
 5x^2-9x-2=0,\quad D=121,\quad x_1=2,\quad x_2=-\frac15.
\]
Первый корень соответствует точке \(A\); ответ: \(\boxed{x_B=-0{,}2}\).

\section*{03. Текстовые задачи на движение}
\ruleline
\begin{checklist}
\checkitem{Прежде всего определяю искомую величину и по возможности обозначаю именно её через \(x\).}
\checkitem{Строю таблицу \(\text{скорость}\;v\mid\text{время}\;t\mid\text{путь}\;S\).}
\checkitem{Один столбец заполняю по формулам \(S=vt\), \(t=S/v\), \(v=S/t\). Это \textbf{формульный столбец}.}
\checkitem{Уравнение составляю по связи величин из условия, используя формульный столбец.}
\checkitem{Согласую единицы: при \(v\) в км/ч путь беру в км, время — в часах.}
\checkitem{Проверяю смысл результата и перевожу ответ в требуемые единицы.}
\end{checklist}

\begin{center}
\renewcommand{\arraystretch}{1.4}
\begin{tabularx}{0.92\linewidth}{>{\bfseries}lXXX}
\toprule
Участник & Скорость \(v\) & Время \(t\) & Путь \(S\)\\
\midrule
Первый & \(v_1\) & \(t_1\) & \(v_1t_1\)\\
Второй & \(v_2\) & \(t_2\) & \(v_2t_2\)\\
\bottomrule
\end{tabularx}
\end{center}

\Needspace{12\baselineskip}
\subsection*{Три связи, которые нужно узнавать}
\begin{enumerate}[label=\textcolor{accent}{\arabic*.}]
\item \key{Навстречу:} до встречи суммы пройденных путей равны расстоянию между пунктами: \(S_1+S_2=L\).
\item \key{В одном направлении:} разность путей равна расстоянию между участниками: \(S_{\text{быстр}}-S_{\text{медл}}=d\) при соответствующем расположении.
\item \key{По кругу:} если быстрый прошёл на один полный круг больше, то \(S_{\text{быстр}}-S_{\text{медл}}=L_{\text{круга}}\).
\end{enumerate}
Если первый выехал на \(\Delta t\) раньше и события сравниваются в один момент, то \(t_{\text{перв}}=t_{\text{втор}}+\Delta t\).

\subsection*{Пример А. Встречное движение с разным временем старта}
Из городов на расстоянии 435 км навстречу едут автомобили со скоростями 60 и 65 км/ч. Первый выехал на час раньше. На каком расстоянии от первого города они встретятся?

Пусть \(x\) км — искомый путь первого. Тогда путь второго \(435-x\), а времени затрачено \(x/60\) и \((435-x)/65\) часов соответственно.
\[
 \frac{x}{60}=\frac{435-x}{65}+1
 \quad\Longrightarrow\quad 65x=60(435-x)+3900
 \quad\Longrightarrow\quad \boxed{x=240\text{ км}}.
\]
Проверка: первый ехал 4 ч, второй — 3 ч, \(240+195=435\).

\subsection*{Пример Б. Движение вдогонку}
Скорости двух пешеходов различаются на 1,5 км/ч. Они вышли одновременно; через какое время расстояние между ними станет 300 м?
\[
 \underbrace{(v+1{,}5)t}_{S_{\text{быстр}}}-\underbrace{vt}_{S_{\text{медл}}}
 =0{,}3\quad\Longrightarrow\quad t=0{,}2\text{ ч}=\boxed{12\text{ мин}}.
\]
\key{Важный приём:} неизвестная \(v\) сократилась; находить её не потребовалось.

\subsection*{Пример В. Один круг разницы}
Длина трассы 14 км. Два автомобиля стартовали одновременно и ехали в одном направлении. Быстрый двигался со скоростью 80 км/ч и за 40 мин прошёл на один круг больше медленного. Найдём скорость \(x\) второго.
\[
 t=\frac{40}{60}=\frac23\text{ ч},\qquad
 80\cdot\frac23-x\cdot\frac23=14
 \quad\Longrightarrow\quad \boxed{x=59\text{ км/ч}}.
\]

\clearpage
\thispagestyle{fancy}
\begin{center}
{\LARGE\bfseries\color{darkaccent}Алгоритм: текстовая задача на движение}\par
\vspace{5mm}
{\small\color{muted}Сначала модель, затем уравнение и ответ}\par
\vspace{6mm}
\end{center}
\tikzset{
 algbox/.style={draw=darkaccent!65,rounded corners=2mm,align=center,text width=8cm,
   minimum height=10mm,inner sep=3mm,font=\small,line width=.65pt,fill=white},
 alga/.style={-Latex,line width=.85pt,draw=accent},
 algchoice/.style={diamond,draw=darkaccent!65,aspect=2.7,align=center,
   text width=3.55cm,inner sep=1mm,font=\small,line width=.65pt},
 algside/.style={algbox,text width=4.2cm,minimum height=14mm}
}
\begin{center}
\begin{tikzpicture}[node distance=9mm]
\node[algbox] (start) {\textbf{Прочитать условие}\newline Что известно и что требуется найти?};
\node[algbox,below=of start] (var) {Обозначить искомое через \(x\)\newline Указать единицы измерения};
\node[algbox,below=of var] (table) {Построить таблицу: \(v\), \(t\), \(S\)\newline Заполнить формульный столбец};
\node[algchoice,below=12mm of table] (decision) {Старт в разное\newline время?};
\node[algside,below left=12mm and 9mm of decision] (yes) {\(t_{\text{раньше}}=t_{\text{позже}}+\Delta t\)};
\node[algside,below right=12mm and 9mm of decision] (no) {Старт одновременно:\newline \(t_1=t_2\)};
\coordinate (merge) at ($(yes.south)!0.5!(no.south)+(0,-12mm)$);
\node[algbox,anchor=north] (relation) at ($(merge)+(0,-5mm)$) {Установить связь путей из текста:\newline сумма при встрече, разность при обгоне};
\node[algbox,below=of relation] (eq) {Составить и решить уравнение};
\node[algbox,below=of eq] (check) {Проверить ограничения и смысл ответа;\newline перевести единицы};
\draw[alga] (start)--(var);
\draw[alga] (var)--(table);
\draw[alga] (table)--(decision);
\draw[alga] (decision.west)-|node[pos=.25,above,sloped,font=\scriptsize]{да}(yes.north);
\draw[alga] (decision.east)-|node[pos=.25,above,sloped,font=\scriptsize]{нет}(no.north);
\draw[draw=accent,line width=.85pt] (yes.south) |- (merge);
\draw[draw=accent,line width=.85pt] (no.south) |- (merge);
\fill[accent] (merge) circle (1.1pt);
\draw[alga] (merge)--(relation.north);
\draw[alga] (relation)--(eq);
\draw[alga] (eq)--(check);
\end{tikzpicture}
\end{center}

\clearpage
\section*{04. Тренировочный блок · Путь А}
\ruleline
\textcolor{muted}{10 заданий, от базовых к составным. Решайте самостоятельно; в прикладных задачах записывайте модель, уравнение или неравенство и ответ с единицами измерения.}
\begin{enumerate}
\item Вычислите: а) $81^{1/4}$; б) $\sqrt[3]{42{,}875}$.
\item Решите уравнение на множестве действительных чисел: $t^4=625$.
\item Решите уравнение: $x^{7/4}=128$.
\item Решите уравнение: $y^{-3/2}={27}$.
\item Вычислите значение выражения $\left({16}\right)^{-3/4}$.
\item На рисунке изображены графики функций $y=\dfrac{k}{x}$ и $y=ax+b$. Найдите абсциссы точек их пересечения.
\begin{center}
\begin{tikzpicture}
\begin{axis}[
    width=0.68\linewidth,
    height=6.8cm,
    axis lines=middle,
    xmin=-3.5, xmax=3.5,
    ymin=-3.5, ymax=4.5,
    samples=300,
    xlabel={$x$}, ylabel={$y$},
    grid=both,
    minor tick num=1,
    clip=false,
    enlargelimits=false,
    ticklabel style={font=\small},
    label style={font=\small},
    every axis plot/.append style={line width=1pt}
]
\addplot[domain=-3.2:-0.25,accent] {2/x};
\addplot[domain=0.25:3.2,accent] {2/x};
\addplot[domain=-3.3:3.3,darkaccent] {x+1};
\addplot[only marks,mark=*,mark size=2pt,textgray] coordinates {(1,2) (-2,-1)};
\node[font=\small,anchor=south west] at (axis cs:1,2) {$A$};
\node[font=\small,anchor=north east] at (axis cs:-2,-1) {$B$};
\end{axis}
\end{tikzpicture}
\end{center}
\item Два автомобиля одновременно стартовали из одной точки круговой трассы длиной 18 км и двигались в одном направлении. Через 45 минут быстрый автомобиль, скорость которого 78 км/ч, прошёл на один круг больше медленного. Найдите скорость медленного автомобиля.
\item По круговой трассе длиной 12 км в одном направлении едут два автомобиля. Медленный автомобиль движется со скоростью 48 км/ч и стартовал первым. Быстрый автомобиль стартовал на 15 минут позже, и спустя 45 минут после своего старта его путь стал на один полный круг больше пути медленного. Найдите скорость быстрого автомобиля.
\item Высота над землёй подброшенного вверх мяча меняется по закону $h(t)=1{,}6+8t-5t^2$, где $h$ --- высота в метрах, $t$ --- время в секундах, прошедшее с момента броска. Сколько секунд мяч будет находиться на высоте не менее трёх метров?
\item После дождя уровень воды в колодце может повыситься. Мальчик измеряет время $t$ падения небольших камешков в колодец и рассчитывает расстояние до воды по формуле $h=5t^2$, где $h$ --- расстояние в метрах, $t$ --- время падения в секундах. До дождя время падения камешков составляло $0{,}6$ с. На сколько должен подняться уровень воды после дождя, чтобы измеряемое время \textbf{уменьшилось} на $0{,}2$ с? Ответ выразите в метрах.
\end{enumerate}

\end{document}
